% ===== 算法与建模可行性分析 =====
% 放在核心公式汇总之后。简单标准方法可直接跳过此节。
% 聚焦最复杂的 1-2 个算法/模型，2-4 个自然段即可。

\subsection{算法与建模可行性分析}

% 分析维度（按需选取，无关的不写）：
%   - 计算可行性：变量/约束规模、求解器能力范围、降维/分解策略、预估求解时间
%   - 收敛性分析：迭代算法收敛条件、收敛速度（线性/超线性/二次）、初值敏感性
%   - 数值稳定性：步长选取依据（如 CFL 条件）、条件数估计、误差传播/累积分析
%   - 适定性分析：解的存在性/唯一性简要论证、可行域是否非空

% 正文示例：
% 针对问题X中使用的【算法/模型名称】，从以下方面分析其可行性与可靠性：
%
% 【计算可行性】问题X的决策变量约 N 个、约束约 M 条，属于【LP/MILP/NLP】问题。
% 采用【求解器/算法】在【参数设置】下求解，单次求解时间约 T 秒，在竞赛资源限制内可行。
%
% 【收敛性】该算法在【条件】下保证收敛，收敛速度为【线性/超线性】。
% 【参数A】选取 B 是为了满足【稳定性条件】。在最坏情况下，误差累积不超过【界】。
%
% 【适用范围】该模型成立的前提是【假设X】。当【条件不满足】时，模型将【退化/失效】。
% 在实际数据范围内（【数据范围】），上述条件均成立。
